% Môn: Vật lý
% Lớp: 10
% Chương: Chương I. Mở đầu
% Bài: L10C1 Bài 2. Các quy tắc an toàn trong phòng thực hành Vật lí · Phân tích nguy cơ mất an toàn liên quan đến thiết bị nhiệt và thủy tinh.
% Số câu: 185
% App đọc file này trực tiếp — sửa rồi Commit trên GitHub, bấm Đồng bộ GitHub .tex

% L10C1 Bài 2. Các quy tắc an toàn trong phòng thực hành Vật lí



% ===== Câu 115 =====
% ID: L10C1B2-34-DS
% Mức: NB
\dangbt{Phân tích nguy cơ mất an toàn liên quan đến thiết bị nhiệt và thủy tinh.}
\begin{ex}
% ID: L10C1B2-34-DS
% Mức: NB
Khi sử dụng bếp điện, cần lưu ý:
\choiceTF
{\True Không chạm tay vào bếp khi bếp đang hoạt động}
{\True Dùng khăn ướt để lau bếp khi bếp đang nóng}
{\True Đặt bếp ở nơi thoáng mát, cách xa vật dễ cháy}
{\True Để thực phẩm gần bếp để tiện nấu nướng}
\loigiai{
\begin{itemchoice}
\item (Đúng) Không chạm tay vào bếp khi bếp đang hoạt động. (Vì): Không chạm tay vào bếp khi đang hoạt động để tránh bị bỏng do nhiệt độ cao
\item (Đúng) Dùng khăn ướt để lau bếp khi bếp đang nóng. (Vì): Lau bếp bằng khăn ướt khi bếp đang nóng có thể gây ra hiện tượng sốc nhiệt và hỏng bếp, cũng như gây bỏng
\item (Đúng) Đặt bếp ở nơi thoáng mát, cách xa vật dễ cháy. (Vì): ặt bếp ở nơi thoáng mát và tránh xa vật dễ cháy để đảm bảo an toàn, phòng tránh nguy cơ cháy nổ
\item (Đúng) Để thực phẩm gần bếp để tiện nấu nướng. (Vì): ể thực phẩm quá gần bếp khi nấu có thể gây cháy, đổ vỡ hoặc mất an toàn
\end{itemchoice}
}
\end{ex}

% ===== Câu 118 =====
% ID: L10C1B2-36-DS
% Mức: TH
\begin{ex}
% ID: L10C1B2-36-DS
% Mức: TH
Khi thực hiện thí nghiệm vật lí, việc tuân thủ các quy tắc an toàn là vô cùng quan trọng. Hãy đánh giá tính đúng/sai của các phát biểu sau đây về các nguy cơ mất an toàn và biện pháp phòng tránh.
\choiceTF
{\True Luôn kiểm tra dây dẫn điện có bị hở hoặc đứt trước khi cắm vào nguồn điện để tránh nguy cơ chập mạch hoặc giật điện.}
{Để tiết kiệm thời gian, có thể bỏ qua việc đọc kỹ hướng dẫn sử dụng thiết bị nếu đã từng làm thí nghiệm tương tự trước đó.}
{\True Khi sử dụng bếp điện hoặc đèn cồn, cần đảm bảo không có vật liệu dễ cháy gần khu vực thí nghiệm để phòng tránh nguy cơ hỏa hoạn.}
{Trong trường hợp xảy ra sự cố nhỏ như đổ hóa chất không độc hại, học sinh nên tự mình xử lý mà không cần báo cáo giáo viên.}
\loigiai{
\begin{itemchoice}
\item (Đúng) Luôn kiểm tra dây dẫn điện có bị hở hoặc đứt trước khi cắm vào nguồn điện để tránh nguy cơ chập mạch hoặc giật điện. (Vì): Việc kiểm tra dây dẫn điện là bước cơ bản và quan trọng để phòng tránh các nguy cơ về điện như chập mạch, đoản mạch hoặc giật điện, đảm bảo an toàn cho người và thiết bị
\item (Sai) Để tiết kiệm thời gian, có thể bỏ qua việc đọc kỹ hướng dẫn sử dụng thiết bị nếu đã từng làm thí nghiệm tương tự trước đó. (Vì): Mỗi thí nghiệm, dù có vẻ tương tự, đều có thể có những yêu cầu, thiết bị hoặc quy trình riêng biệt. Việc bỏ qua hướng dẫn sử dụng có thể dẫn đến việc vận hành sai, làm hỏng thiết bị hoặc gây nguy hiểm cho người thực hiện
\item (Đúng) Khi sử dụng bếp điện hoặc đèn cồn, cần đảm bảo không có vật liệu dễ cháy gần khu vực thí nghiệm để phòng tránh nguy cơ hỏa hoạn. (Vì): Bếp điện và đèn cồn là nguồn nhiệt có thể gây cháy. Việc giữ khoảng cách an toàn với vật liệu dễ cháy là quy tắc cơ bản để phòng tránh hỏa hoạn trong phòng thí nghiệm
\item (Sai) Trong trường hợp xảy ra sự cố nhỏ như đổ hóa chất không độc hại, học sinh nên tự mình xử lý mà không cần báo cáo giáo viên. (Vì): Mọi sự cố trong phòng thí nghiệm, dù lớn hay nhỏ, đều cần được báo cáo cho giáo viên hoặc người phụ trách. Điều này giúp đảm bảo sự cố được xử lý đúng cách, ngăn ngừa các rủi ro tiềm ẩn và rút kinh nghiệm cho các lần sau
\end{itemchoice}
}
\end{ex}

% ===== Câu 167 =====
% ID: L10C1B2-75-TN
% Mức: NB
\begin{ex}
% ID: L10C1B2-75-TN
% Mức: NB
Khi một ống nghiệm thủy tinh bị vỡ trong phòng thực hành, học sinh nên làm gì?
\choice
{Dùng tay nhặt ngay các mảnh vỡ}
{Dùng chân gạt các mảnh vỡ vào góc bàn}
{\True Báo giáo viên và dùng dụng cụ phù hợp để thu gom}
{Tiếp tục làm thí nghiệm như bình thường}
\loigiai{
Mảnh thủy tinh vỡ sắc nhọn có thể gây đứt tay. Cần báo giáo viên và thu gom bằng dụng cụ phù hợp.
}
\end{ex}

% ===== Câu 168 =====
% ID: L10C1B2-78-TN
% Mức: TH
\begin{ex}
% ID: L10C1B2-78-TN
% Mức: TH
Vì sao không nên rót nước lạnh đột ngột vào bình thủy tinh đang nóng?
\choice
{Vì nước lạnh làm tăng khối lượng bình}
{\True Vì bình có thể nứt vỡ do thay đổi nhiệt độ đột ngột}
{Vì nước lạnh làm bình nhiễm điện}
{Vì nước lạnh làm mất màu thủy tinh}
\loigiai{
Sự thay đổi nhiệt độ đột ngột làm các phần của thủy tinh co giãn không đều, có thể gây nứt vỡ.
}
\end{ex}

% ===== Câu 169 =====
% ID: L10C1B2-79-TN
% Mức: TH
\begin{ex}
% ID: L10C1B2-79-TN
% Mức: TH
Bình thủy tinh chịu nhiệt vừa đun nóng ở nhiệt độ cao nếu bị dội ngay nước lạnh vào thì có nguy cơ xảy ra hiện tượng gì?
\choice
{Bình thủy tinh bị co lại làm tăng độ dày của thành bình}
{\True Bình thủy tinh bị nứt vỡ đột ngột do sự co giãn vì nhiệt diễn ra không đều giữa các phần}
{Nước lạnh làm tăng độ bền cơ học giúp bình chịu va đập tốt hơn}
{Không xảy ra hiện tượng gì vì bình thủy tinh chịu nhiệt có hệ số nở dài bằng không}
\loigiai{
Sự thay đổi nhiệt độ quá đột ngột khiến lớp thủy tinh bên ngoài co lại nhanh chóng trong khi lớp bên trong chưa kịp co theo, tạo ra ứng suất cơ học rất lớn làm nứt vỡ bình thủy tinh.
}
\end{ex}

% ===== Câu 170 =====
% ID: L10C1B2-80-TN
% Mức: TH
\begin{ex}
% ID: L10C1B2-80-TN
% Mức: TH
Tại sao học sinh không được dùng tay trần chạm trực tiếp lên bề mặt làm việc của thấu kính và gương phẳng trong phòng thực hành quang hình?
\choice
{Vì bề mặt thấu kính rất sắc nhọn có thể làm đứt tay học sinh}
{\True Vì mồ hôi và vân tay bám vào sẽ gây mốc, trầy xước và làm sai lệch đường truyền của các tia sáng}
{Vì thấu kính có thể tích điện truyền dòng điện tĩnh sang cơ thể người gây giật nhẹ}
{Vì tay trần tiếp xúc sẽ làm thay đổi chiết suất của chất liệu làm thấu kính vĩnh viễn}
\loigiai{
Bề mặt quang học của thấu kính và gương phẳng rất nhạy cảm. Vân tay và mồ hôi chứa axit yếu và ẩm mốc sẽ ăn mòn lớp phủ quang học, gây tán xạ ánh sáng và làm sai lệch nghiêm trọng kết quả thí nghiệm.
}
\end{ex}

% ===== Câu 171 =====
% ID: L10C1B2-82-TN
% Mức: TH
\begin{ex}
% ID: L10C1B2-82-TN
% Mức: TH
Khi đang tiến hành thí nghiệm đo nhiệt độ bằng nhiệt kế thủy ngân, không may nhiệt kế bị rơi vỡ và thủy ngân tràn ra sàn phòng thực hành. Học sinh cần thực hiện hành động nào sau đây để đảm bảo an toàn?
\choice
{Dùng tay không hoặc khăn giấy để thu gom các mảnh vỡ và giọt thủy ngân.}
{\True Báo cáo ngay cho giáo viên hoặc người phụ trách phòng thí nghiệm và không chạm vào chất lỏng thủy ngân.}
{Mở cửa sổ để thông gió và tiếp tục thí nghiệm với nhiệt kế khác.}
{Dùng chổi quét sạch các mảnh vỡ và thủy ngân vào thùng rác thông thường.}
\loigiai{
Thủy ngân là một kim loại lỏng độc hại, có thể gây nguy hiểm khi tiếp xúc trực tiếp hoặc hít phải hơi của nó. Khi nhiệt kế thủy ngân bị vỡ, hành động đúng đắn và an toàn nhất là không được chạm vào thủy ngân, mà phải báo cáo ngay cho giáo viên hoặc người phụ trách phòng thí nghiệm để họ có biện pháp xử lý chuyên nghiệp, đảm bảo an toàn cho mọi người và môi trường. Các phương án $A$, $C$, $D$ đều tiềm ẩn nguy cơ gây hại cho sức khỏe do tiếp xúc trực tiếp với thủy ngân hoặc phát tán hơi thủy ngân độc hại.
}
\end{ex}

% ===== Câu 172 =====
% ID: L10C1B2-83-TN
% Mức: TH
\begin{ex}
% ID: L10C1B2-83-TN
% Mức: TH
Trong phòng thực hành Vật lí, khi tiến hành thí nghiệm với mạch điện có nguồn điện áp cao (trên $40\,\text{V}$), một học sinh sơ ý chạm tay trực tiếp vào các điểm nối hở của mạch. Nguy cơ mất an toàn nghiêm trọng nhất mà học sinh này có thể gặp phải là gì?
\choice
{\True Bị điện giật gây co giật cơ, bỏng hoặc ngừng tim.}
{Làm hỏng các linh kiện điện tử trong mạch do đoản mạch.}
{Gây ra hiện tượng quá tải, dẫn đến cháy nổ thiết bị.}
{Làm rơi vỡ các dụng cụ thủy tinh đặt gần khu vực thí nghiệm.}
\loigiai{
Khi làm việc với mạch điện áp cao, việc chạm tay trực tiếp vào các điểm nối hở của mạch là cực kỳ nguy hiểm. Dòng điện có thể đi qua cơ thể gây ra điện giật, dẫn đến các hậu quả nghiêm trọng như co giật cơ, bỏng nặng, rối loạn nhịp tim hoặc thậm chí ngừng tim, đe dọa trực tiếp đến tính mạng. Các nguy cơ khác như hỏng thiết bị hay cháy nổ cũng có thể xảy ra nhưng nguy cơ điện giật trực tiếp đến người là nghiêm trọng và tức thời nhất trong tình huống này.
}
\end{ex}

% ===== Câu 173 =====
% ID: L10C1B2-84-TN
% Mức: TH
\begin{ex}
% ID: L10C1B2-84-TN
% Mức: TH
Trong một buổi thực hành thí nghiệm về đo gia tốc trọng trường, một học sinh làm rơi một vật nhỏ từ độ cao nhất định. Theo quy tắc an toàn trong phòng thực hành Vật lí, hành động nào sau đây là KHÔNG phù hợp?
\choice
{Đảm bảo khu vực xung quanh nơi vật rơi không có người hoặc vật cản.}
{Sử dụng kính bảo hộ để tránh vật rơi hoặc mảnh vỡ bắn vào mắt.}
{\True Thả vật rơi từ vị trí cao nhất có thể để quan sát rõ quỹ đạo.}
{Kiểm tra độ chắc chắn của giá đỡ hoặc thiết bị giữ vật trước khi thả.}
\loigiai{
Trong phòng thực hành, việc đảm bảo an toàn là ưu tiên hàng đầu. Thả vật rơi từ vị trí cao nhất có thể tiềm ẩn nguy cơ gây hại cho người xung quanh hoặc làm hỏng thiết bị nếu vật rơi không đúng vị trí hoặc không kiểm soát được. Các hành động $A$, $B$, $D$ đều là những biện pháp an toàn cần thiết khi thực hiện thí nghiệm có liên quan đến vật rơi.
}
\end{ex}

